% Esta sección presenta la evaluación experimental de la estrategia de almacenamiento comprimido
% integrada en el simulador. El objetivo no es únicamente reportar diferencias de desempeño,
% sino establecer bajo qué condiciones la compresión del vector de estado constituye una
% alternativa conveniente frente al almacenamiento no comprimido. Para ello se comparan tres
% estrategias de representación: una referencia sin compresión, una propuesta sin pérdida
% basada en Blosc y una estrategia con pérdida controlada basada en ZFP.

% La evaluación se organizó sobre dos familias de circuitos con comportamientos contrastantes
% desde el punto de vista de la compresibilidad: el algoritmo de búsqueda de Grover y la
% transformada cuántica de Fourier (QFT). Esta selección permite observar la propuesta tanto en
% escenarios donde el estado cuántico conserva regularidades explotables como en otros donde
% esa regularidad es menor y, por tanto, la compresión enfrenta condiciones menos favorables.

% La campaña experimental cubrió 22, 23 y 24 qubits. En Grover se estudiaron dos variantes:
% objetivo único y multiobjetivo. En QFT se consideraron dos tipos de entrada: una
% superposición periódica y una superposición de alta entropía con fases aleatorias. Con esta
% combinación se buscó evitar una lectura sesgada por un solo patrón de estado y, al mismo
% tiempo, contrastar comportamientos representativos de compresibilidad alta y baja.

% Las tablas de resultados reportan cuatro magnitudes principales: memoria residente máxima,
% ahorro de memoria respecto a la referencia no comprimida, tiempo total de ejecución y tiempo
% relativo expresado en múltiplos de la referencia. En el caso de Blosc, al tratarse de una
% estrategia sin pérdida, la exactitud se verificó comparando amplitud por amplitud el vector
% completo exportado por cada ejecución comprimida frente a la ejecución no comprimida de la
% misma instancia. En todos los casos reportados, esa comparación produjo igualdad exacta y,
% por tanto, error máximo absoluto nulo. Para ZFP, la misma comparación completa se utilizó
% para calcular el máximo error absoluto frente a la referencia.

% Además de esta comparación principal, en Grover se incorpora una evaluación complementaria
% con la variante de parametrización reducida descrita en la
% Sección~\ref{subsec:grover_reducido_impl}. Su propósito es distinguir entre el beneficio que
% proviene de comprimir el almacenamiento del estado y el que resulta de explotar directamente
% la estructura algebraica del algoritmo. En esa subsección complementaria, por tanto, la
% lectura se concentra en memoria y tiempo, ya que el interés principal no es contrastar
% exactitud entre compresores, sino delimitar el papel relativo de cada estrategia de
% optimización.
